\chapter[Variational Structure]{Variational Structure and Interior Critical Points of the Uniform Observer Divergence}
\section{Introduction}
The preceding chapters established the fundamental analytical properties of the Uniform Observer Divergence (UOD), including its interpretation as the variance of the log-likelihood ratio under the uniform observer, its invariance properties, and its relationship with classical information-theoretic divergences.
In this chapter we investigate a different aspect of the functional, namely its variational geometry.
Rather than evaluating the UOD for a prescribed pair of probability distributions, we consider the inverse problem: among all probability distributions having a fixed Kullback--Leibler divergence from the uniform distribution, which configurations maximize the disagreement measured by the UOD?
More precisely, throughout this chapter we study the constrained optimization problem
\[
\max_{P\in\Delta_n^\circ}
\mathcal D(P,U_n)
\]
subject to
$D_{\mathrm{KL}}(P\|U_n)=K,$
where \(U_n\) denotes the uniform distribution on \(n\) symbols, \(K>0\) is fixed, and \(\Delta_n^\circ\) is the interior of the probability simplex.
At first sight this appears to be a nonlinear optimization problem in \(n\) independent variables. One might therefore expect the geometry of its stationary points to become increasingly complicated as the ambient dimension grows.
The main result of this chapter shows that the opposite is true.
Although the optimization is performed over an \(n\)-dimensional manifold, every interior stationary point necessarily collapses to a configuration containing at most three distinct log-likelihood values.
This dimensional collapse is completely independent of the ambient dimension and follows from a remarkable interaction between the stationarity equations and the classical theory of Extended Chebyshev Systems.
Throughout this chapter we restrict our attention to \textbf{interior regular constrained critical points}. The characterization of global extrema, which may occur on the boundary of the simplex, is left for future work.
\section{Variational Formulation}
Throughout this chapter the reference distribution is fixed to be the uniform distribution \(Q=U_n\).
Let \(P=(P_1,\ldots,P_n)\in\Delta_n^\circ\). Since every coordinate is strictly positive, introduce the log-likelihood coordinates
\[
a_i
=
\log\frac{P_i}{U_i}
=
\log(nP_i),\qquad
i=1,\ldots,n.
\]
The inverse transformation is \(P_i=e^{a_i}/n\), which automatically guarantees positivity.
Normalization of the probability distribution becomes
$\sum_{i=1}^{n}e^{a_i}=n.$
The KL divergence constraint becomes
\[
D_{\mathrm{KL}}(P\|U_n)
=
\frac1n\sum_i e^{a_i}a_i
=
K.
\]
The objective functional is the UOD with respect to the uniform distribution,
\[
\mathcal D(P,U_n)
=
\frac1n\sum_i a_i^2
-
\left(\frac1n\sum_i a_i\right)^2
=
\operatorname{Var}_U(a).
\]
Thus, in log-likelihood coordinates, the optimization problem becomes
\[
\max_{a\in\mathbb R^n}
\frac1n\sum_i a_i^2
-
\left(\frac1n\sum_i a_i\right)^2
\]
subject to
\[
\frac1n\sum_i e^{a_i}=1,
\qquad
\frac1n\sum_i e^{a_i}a_i=K.
\]
\section{Regularity of the Constraint Manifold}
\begin{lemma}[Regularity of the Constraint Manifold]
\label{ch:lem:29-1}
Assume \(K>0\). Then every feasible point satisfies the Linear Independence Constraint Qualification (LICQ).
\end{lemma}
\begin{proof}
The gradients of the two constraints are
\[
\nabla g_1
=
\frac1n
\begin{pmatrix}
e^{a_1}, \ldots, e^{a_n}
\end{pmatrix},
\qquad
\nabla g_2
=
\frac1n
\begin{pmatrix}
e^{a_1}(1+a_1), \ldots, e^{a_n}(1+a_n)
\end{pmatrix}.
\]
These two vectors are linearly dependent only if there exists a scalar \(\lambda\) such that \(e^{a_i}(1+a_i)=\lambda e^{a_i}\) for all \(i\), which would imply all coordinates \(a_i\) are equal. This forces \(P=U_n\) and consequently \(K=0\), contradicting \(K>0\). Therefore the gradients are linearly independent everywhere on the constraint manifold, establishing LICQ. 
\end{proof}
%% Explanatory comment: The LICQ condition guarantees that the constraint manifold of feasible probability distributions is regular, meaning that the Lagrange multiplier method is valid for all interior critical points. Consequently, every constrained stationary point satisfies the Euler–Lagrange equations derived in the next section.
\section{Euler--Lagrange Equations}
By Lemma~\ref{ch:lem:29-1}, every interior constrained critical point satisfies the Euler--Lagrange equations
$\nabla\mathcal D = \lambda_1\nabla g_1 + \lambda_2\nabla g_2$
for some Lagrange multipliers \(\lambda_1,\lambda_2\in\mathbb R\).
Computing the gradients,
\[
\frac{\partial\mathcal D}{\partial a_i}
=
\frac{2}{n}(a_i-\bar a),
\]
\[
\frac{\partial g_1}{\partial a_i}
=
\frac{1}{n}e^{a_i},
\qquad
\frac{\partial g_2}{\partial a_i}
=
\frac{1}{n}e^{a_i}(1+a_i).
\]
Thus the stationarity condition becomes
\[
\frac{2}{n}(a_i-\bar a)
=
\frac{\lambda_1}{n}e^{a_i}
+
\frac{\lambda_2}{n}e^{a_i}(1+a_i),
\]
which simplifies to
\[
2(a_i-\bar a) = e^{a_i}\bigl(\lambda_1 + \lambda_2(1+a_i)\bigr),\qquad i=1,\ldots,n.
\]
Let \(\alpha = \lambda_1+\lambda_2\) and \(\beta = \lambda_2\). Then
$2(a_i-\bar a) = e^{a_i}(\alpha + \beta a_i).$
\section{The Stationarity Function}
Define the function
$\varphi(x) = 2(x-\bar a) - e^{x}(\alpha + \beta x).$
Then the stationarity equations are simply \(\varphi(a_i)=0\) for all \(i\).
Thus every coordinate of a critical point must be a root of the same scalar function \(\varphi\), which depends on the three parameters \(\bar a,\alpha,\beta\) determined by the constraints.
\begin{lemma}[Structure of the Stationarity Function]
\label{ch:lem:29-2}
For any real parameters \(\bar a,\alpha,\beta\) with \(\beta\neq0\), the function \(\varphi\) is real-analytic and has at most three distinct real zeros.
\end{lemma}
\begin{proof}
Consider the auxiliary function
$\psi(x)=e^{-x}\varphi(x)=2e^{-x}(x-\bar a)-(\alpha+\beta x).$
Differentiating,
$\psi'(x)=2e^{-x}(1-x+\bar a)-\beta.$
Since \(e^{-x}(1-x+\bar a)\) is bounded and \(\beta\neq0\), the equation \(\psi'(x)=0\) has finitely many solutions. By the theory of exponential polynomials, \(\psi\) can have at most three zeros unless it vanishes identically. The identical vanishing case would imply \(\varphi\equiv0\), which is impossible for nontrivial parameters because the exponential term cannot be canceled by the linear term globally. Therefore \(\varphi\) has at most three distinct real zeros. 
\end{proof}
%% Explanatory comment: This lemma identifies the key structural property of the stationarity function: regardless of the ambient dimension n, every coordinate of a critical point solves the same scalar equation, which admits at most three solutions. This observation is the first step toward the dimensional collapse phenomenon.
\begin{corollary}[Coordinate Collapse]
\label{ch:cor:29-1}
Every interior regular constrained critical point has at most three distinct values among its coordinates \(a_1,\ldots,a_n\).
\end{corollary}
\begin{proof}
Since every coordinate \(a_i\) satisfies \(\varphi(a_i)=0\) and \(\varphi\) has at most three distinct real zeros (Lemma~\ref{ch:lem:29-2}), the coordinates can take at most three distinct values. 
\end{proof}
%% Explanatory comment: The corollary of the coordinate collapse is a striking consequence: every interior constrained critical point concentrates into at most three distinct log-likelihood values. This reduction from an n-dimensional coordinate vector to at most three distinct values is the central phenomenon of the chapter.
\section{Extended Chebyshev Structure}
\begin{lemma}[Extended Chebyshev System]
\label{ch:lem:29-3}
The set of functions \(\{1,\;x,\;e^x,\;xe^x\}\) forms an Extended Chebyshev System on the real line.
\end{lemma}
\begin{proof}
The successive principal Wronskians are
\[
W_1(x)=1\neq0,\quad
W_2(x)=
\begin{vmatrix}
1 & x\\ 0 & 1
\end{vmatrix}=1\neq0,
\]
\[
W_3(x)=
\begin{vmatrix}
1 & x & e^x\\
0 & 1 & e^x\\
0 & 0 & e^x
\end{vmatrix}=e^x\neq0,
\]
\[
W_4(x)=
\begin{vmatrix}
1 & x & e^x & xe^x\\
0 & 1 & e^x & e^x(1+x)\\
0 & 0 & e^x & e^x(2+x)\\
0 & 0 & e^x & e^x(3+x)
\end{vmatrix}=e^{3x}\neq0.
\]
All Wronskians are non-vanishing, confirming the Extended Chebyshev property. 
\end{proof}
%% Explanatory comment: The set of functions {1, x, e^x, xe^x} forms an Extended Chebyshev System, which is the deepest structural explanation for the three-level collapse. This property implies that any nontrivial linear combination of these functions has at most three zeros, directly bounding the number of distinct coordinate values at a critical point.
\begin{lemma}[Non-Degeneracy]
\label{ch:lem:29-4}
The stationarity function \(\varphi\) cannot vanish identically on any interval of positive length.
\end{lemma}
\begin{proof}
Identical vanishing would imply \(2(x-\bar a)=e^x(\alpha+\beta x)\) for all \(x\) in an interval. If \(\beta=0\), then \(2(x-\bar a)=\alpha e^x\), which cannot hold identically because the left-hand side is linear while the right-hand side is exponential. If \(\beta\neq0\), the exponential term grows faster than the linear term, preventing identical equality. Therefore \(\varphi\) cannot vanish identically. 
\end{proof}
%% Explanatory comment: The stationarity function cannot vanish identically on any interval, confirming that the root-counting bound from the Chebyshev System is not artificially saturated. This non-degeneracy ensures that the three-level structure is a genuine geometric constraint rather than a trivial consequence of a vanishing function.
\begin{lemma}[Zero Counting]
\label{ch:lem:29-5}
The function \(\varphi\) with \(\beta\neq0\) has at most three distinct real zeros. Moreover, if it has exactly three distinct zeros, they are simple and the sign pattern of \(\varphi\) alternates between them.
\end{lemma}
\begin{proof}
The function \(\varphi\) belongs to the four-dimensional space spanned by \(\{1,x,e^x,xe^x\}\). By Lemma~\ref{ch:lem:29-3}, this space is an Extended Chebyshev System. By Lemma~\ref{ch:lem:29-4}, \(\varphi\) is not identically zero. A fundamental property of Extended Chebyshev Systems is that any non-zero element has at most three zeros. The simplicity and sign alternation follow from the non-vanishing of the Wronskians. 
\end{proof}
%% Explanatory comment: This lemma sharpens the bound: the stationarity function has at most three zeros, and when exactly three occur they are simple with alternating sign patterns. The simplicity property is important because it guarantees that the critical points are non-degenerate, which in turn implies structural stability under small perturbations.
\section{Main Structural Theorem}
\begin{theorem}[Three-Level Structure of Interior Regular Constrained Critical Points]
\label{ch:thm:29-3}
Every interior regular constrained critical point of the maximization problem
\[
\max_{P\in\Delta_n^\circ}\mathcal D(P,U_n)\quad\text{subject to}\quad D_{\mathrm{KL}}(P\|U_n)=K>0
\]
has at most three distinct log-likelihood values. Consequently, the corresponding probability distribution has at most three distinct probability values.
Moreover, if exactly three distinct values occur, they are arranged as \(a_{i_1}<a_{i_2}<a_{i_3}\) with multiplicities \(m_1,m_2,m_3\ge1\), and the stationarity equations reduce to a finite-dimensional system of five equations in five unknowns, together with the two constraint equations.
\end{theorem}
\begin{proof}
By Lemma~\ref{ch:lem:29-1}, LICQ holds, so the Euler--Lagrange equations are necessary. The stationarity equations give \(\varphi(a_i)=0\) for every coordinate \(i\). By Lemma~\ref{ch:lem:29-5}, \(\varphi\) has at most three distinct real zeros. Hence the coordinates take at most three distinct values.
The probability values are obtained through the transformation \(P_i=e^{a_i}/n\), which preserves distinctness.
When exactly three distinct values occur, let them be \(a_1<a_2<a_3\) with multiplicities \(m_1,m_2,m_3\ge1\) and \(m_1+m_2+m_3=n\). The unknowns are: the three coordinates \(a_1,a_2,a_3\), the three multiplicities \(m_1,m_2,m_3\), and the two Lagrange multipliers \(\lambda_1,\lambda_2\) (equivalently \(\alpha,\beta\)). These satisfy: the two constraint equations (normalization and KL divergence), the three stationarity equations \(\varphi(a_j)=0\) for \(j=1,2,3\), and the multiplicity sum \(m_1+m_2+m_3=n\). This yields a finite system determining the possible configurations. 
\end{proof}
%% Explanatory comment: The Three-Level Structure Theorem is the central result of this chapter: every interior regular constrained critical point collapses to at most three distinct log-likelihood values, regardless of the ambient dimension n. This dimensional collapse is universal and follows from the interaction between the stationarity equations and the Extended Chebyshev System structure.
\begin{corollary}[Finite-Dimensional Reduction]
\label{ch:cor:29-2}
The search for interior regular constrained critical points reduces from the infinite-dimensional space of probability distributions on \(n\) symbols to a finite-dimensional algebraic system of at most five equations in five unknowns, independently of \(n\).
\end{corollary}
\begin{proof}
Theorem~\ref{ch:thm:29-3} explicitly identifies the unknowns and the governing equations. The reduction is independent of \(n\) because the structure of the stationarity function \(\varphi\) depends only on the parameters \(\bar a,\alpha,\beta\) and not on the dimension. 
\end{proof}
%% Explanatory comment: The finite-dimensional reduction transforms what appears to be a high-dimensional optimization problem into a small algebraic system of at most five equations in five unknowns, independently of n. This is the practical payoff of the three-level structure: the search for interior critical points becomes computationally tractable even in very high dimensions.
\section{Explicit Characterization for Three Distinct Values}
When three distinct log-likelihood values \(a<b<c\) occur with multiplicities \(p,q,r\ge1\) and \(p+q+r=n\), the system becomes:
Normalization: \(p e^{a}+q e^{b}+r e^{c}=n\).
KL constraint: \(p e^{a}a+q e^{b}b+r e^{c}c=nK\).
Stationarity at each distinct value:
\[
2(a-\bar a)=e^{a}(\alpha+\beta a),\quad
2(b-\bar a)=e^{b}(\alpha+\beta b),\quad
2(c-\bar a)=e^{c}(\alpha+\beta c),
\]
where \(\bar a=(p a+q b+r c)/n\).
This is a system of five equations in the unknowns \(a,b,c,\alpha,\beta\) plus the integer multiplicities \(p,q,r\).

\begin{theorem}[Explicit Parametrization of Three-Level Critical Points]
\label{ch:thm:three-level-param}
Let \(P\) be an interior regular constrained critical point with three distinct log-likelihood values \(a<b<c\), multiplicities \(p,q,r\ge1\) with \(p+q+r=n\), and Lagrange multipliers \(\alpha,\beta\) (equivalently \(\lambda_1,\lambda_2\)).  Then the three values are determined by the two parameters \(\alpha,\beta\) through the equation

\[
2(a-\bar a) = e^{a}(\alpha+\beta a),\qquad
2(b-\bar a) = e^{b}(\alpha+\beta b),\qquad
2(c-\bar a) = e^{c}(\alpha+\beta c).
\]

Subtracting the first two equations eliminates \(\bar a\) and \(\alpha\):

\[
2(b-a) = e^{b}(\alpha+\beta b) - e^{a}(\alpha+\beta a).
\]

Solving for \(\alpha\) in terms of \(a,b,\beta\) gives

\[
\alpha = \frac{2(b-a) - \beta(b e^{b} - a e^{a})}{e^{b} - e^{a}}.
\]

Substituting back, the third equation becomes a single transcendental
equation in the three variables \(a,b,c\) with parameter \(\beta\):

\[
\frac{2(c-a) - \beta(c e^{c} - a e^{a})}{e^{c} - e^{a}}
=
\frac{2(b-a) - \beta(b e^{b} - a e^{a})}{e^{b} - e^{a}}.
\]

For fixed \(\beta\), this equation determines a two-parameter family of
triples \((a,b,c)\).  The constraints (normalization and KL divergence)
then determine the multiplicities \(p,q,r\) uniquely (when they exist as
positive integers summing to \(n\)).
\end{theorem}

\begin{corollary}[Two-Parameter Family]
\label{ch:cor:three-level-family}
The set of all interior regular constrained critical points with
exactly three distinct values forms a two-parameter family, parameterized
by \(\beta\neq0\) and one of the log-likelihood values (e.g.\ \(a\)).
The remaining two values \(b,c\) and the multiplicities \(p,q,r\) are
determined by the constraints, provided the resulting multiplicities are
positive integers.
\end{corollary}

\begin{proof}
The system has \(5\) unknowns \((a,b,c,\alpha,\beta)\) and
\(5\) equations (three stationarity equations, normalization, KL
constraint), but one degree of freedom is absorbed by the implicit
relation \(\bar a = (pa+qb+rc)/n\).  Fixing \(\beta\) and \(a\) leaves
\(b,c\) determined by the stationarity equations and the constraints,
generically giving isolated solutions. \(\)
\end{proof}

\begin{example}[Binary Case \(n=2\)]
\label{ch:ex:three-level-binary}
For \(n=2\) and \(K>0\), the Three-Level Theorem allows at most three
values, but with only two coordinates the only possibility is at most
two distinct values.  Let \(a<b\) with multiplicities \(p=1,q=1\).
The stationarity equations reduce to

\[
2(a-\bar a)=e^{a}(\alpha+\beta a),\qquad
2(b-\bar a)=e^{b}(\alpha+\beta b),
\]

with \(\bar a=(a+b)/2\).  Subtracting and solving gives

\[
b-a = \frac{1}{2}\bigl(e^{b}(\alpha+\beta b) - e^{a}(\alpha+\beta a)\bigr).
\]

For a given \(K\), the system determines \((a,b,\alpha,\beta)\) uniquely.
The solution can be verified to be symmetric: \(P=(p,1-p)\) with
\(p>1/2\), and the KL constraint \(K = p\log(2p)+(1-p)\log(2(1-p))\)
determines \(p\).
\end{example}

\begin{example}[Symmetric Three-Level Configuration]
\label{ch:ex:three-level-symmetric}
Consider a symmetric configuration where the three log-likelihood values
are equally spaced: let \(a=-\delta\), \(b=0\), \(c=\delta\) with
\(\delta>0\).  By symmetry, the multiplicities satisfy \(p=r\) and
\(q=n-2p\).  The stationarity equations at \(\pm\delta\) become

\[
2(-\delta) = e^{-\delta}(\alpha-\beta\delta),\qquad
2\delta = e^{\delta}(\alpha+\beta\delta).
\]

Adding the two equations gives
\(0 = e^{-\delta}(\alpha-\beta\delta) + e^{\delta}(\alpha+\beta\delta)\),
which implies \(\alpha(e^{\delta}+e^{-\delta}) = \beta\delta(e^{\delta}-e^{-\delta})\), or

\[
\alpha = \beta\delta\tanh\delta.
\]

Substituting back, the stationarity equation at \(b=0\) gives

\[
2(0-\bar a) = \alpha,
\]

so \(\bar a = -\alpha/2 = -(\beta\delta\tanh\delta)/2\).

The normalization constraint \(p e^{-\delta} + q + p e^{\delta} = n\) and
the KL constraint then determine the admissible values of \(\delta\)
and \(p\) for each \(n\).  For \(n=3\) and \(p=1,q=1\), the system
admits a unique solution, which is the global maximum of the UOD under
the KL constraint for that dimension.
\end{example}

\begin{remark}[Existence and Uniqueness]
\label{rem:three-level-existence}
The existence of interior critical points for a given \(K>0\) is not
guaranteed.  The symmetric example above shows that for small \(K\),
the only possibility is the uniform distribution (the trivial
one-level case).  As \(K\) increases, a pitchfork bifurcation creates
asymmetric two-level configurations, and for sufficiently large \(K\),
three-level configurations appear.  The precise bifurcation values
depend on \(n\) and the multiplicities.
\end{remark}

\section{Discussion: What the Three-Level Structure Implies}

The Three-Level Theorem and the explicit parametrization derived above
have several consequences for the geometry of the UOD and for the
information-theoretic bounds developed in Parts~III--IV of this book.

\paragraph{Dimensional reduction is universal.}
The most striking implication is that the variational problem of
maximising the UOD under a KL constraint, despite being formulated over
an \(n\)-dimensional manifold, is governed by a system of equations
whose complexity does not grow with \(n\).  Whether \(n=3\) or
\(n=10^6\), every interior critical point is described by at most three
distinct probability values and their integer multiplicities.  This is
not an asymptotic approximation but an exact structural collapse.

\paragraph{Tightness of the geometric bounds.}
The configurations that saturate the Lipschitz and quadratic bounds of
Theorems~\ref{ch:thm:14-1} and~\ref{ch:thm:14-2} are themselves
solutions of constrained optimisation problems on the posterior
manifold.  The three-level critical points identified here are the
natural candidates for extremal configurations: a distribution that
maximises the UOD for a given KL divergence from the uniform
distribution is likely to be a three-level configuration (or a
boundary limit thereof).  If such a configuration can be shown to
saturate a bound for a given functional (say, the Lipschitz bound for
the R\'enyi entropy), then that bound is tight.

Conversely, if a bound cannot be saturated by any three-level
configuration, it cannot be saturated at all by an interior point.
This provides a practical, finite-dimensional test for the tightness of
any bound derived from the universal geometric framework.

\paragraph{A pitchfork bifurcation governs the transition from
uniformity.}
The uniform distribution \(U_n\) is a trivial one-level critical point
for \(K=0\).  As \(K\) increases past a threshold \(K_1(n)\), the
uniform solution loses stability and a two-level family appears.
For sufficiently large \(K>K_2(n)\), three-level configurations emerge.
This sequence---one level, two levels, three levels---is a
pitchfork bifurcation of the stationarity equations, and it is the same
qualitative behaviour for every \(n\).  The thresholds depend on \(n\)
only through the integer multiplicities, not through the continuous
geometry.

\paragraph{Relation to the Extended Chebyshev System.}
The dimensional collapse is not an accident of the specific objective
and constraint.  The functions \(1,\;x,\;e^x,\;xe^x\) that appear in the
stationarity equations span a four-dimensional Extended Chebyshev
System, and it is this property---the bounded zero-count of linear
combinations---that forces the coordinate collapse.  Any variational
problem whose Euler--Lagrange equations produce the same function space
will exhibit the same phenomenon.  This suggests that the three-level
structure may appear in other contexts, whenever two exponential-type
constraints interact with a quadratic objective.

\paragraph{Open question: boundary maxima.}
The Three-Level Theorem applies only to interior critical points.
The global maximum of the UOD under the KL constraint may lie on the
boundary of the probability simplex, where some probabilities vanish
and the log-likelihood coordinates are no longer defined.  Such
boundary configurations are not captured by the present analysis.
However, the interior classification provides a necessary condition for
any interior extremum, and the boundary case requires separate
treatment---likely involving a limit analysis of the three-level
configurations as some multiplicities tend to zero.
\subsection{Boundary Behavior}
Theorem~\ref{ch:thm:29-3} characterizes only interior critical points. The global maximum of the UOD under the KL constraint may occur on the boundary of the probability simplex, where some coordinates vanish.
On the boundary, the log-likelihood coordinates are no longer defined, and a separate analysis is required.
However, any interior critical point must satisfy the three-level structure identified above. This provides a complete classification of all interior stationary points, regardless of the ambient dimension.


Theorem~\ref{ch:thm:29-3} characterizes only interior critical points. The global maximum of the UOD under the KL constraint may occur on the boundary of the probability simplex, where some coordinates vanish.
On the boundary, the log-likelihood coordinates are no longer defined, and a separate analysis is required.
However, any interior critical point must satisfy the three-level structure identified above. This provides a complete classification of all interior stationary points, regardless of the ambient dimension.
\section{Summary}
This chapter has established the following results.
\begin{itemize}
\item The problem of maximizing the UOD subject to a fixed KL divergence from the uniform distribution admits a clean variational formulation in log-likelihood coordinates, with a quadratic objective and exponential constraints (Section 29.2).
\item The constraint manifold is regular (LICQ holds) for \(K>0\) (Section 29.3).
\item Every interior regular constrained critical point satisfies a scalar stationarity equation \(\varphi(a_i)=0\) for each coordinate (Section 29.4).
\item The stationarity function \(\varphi\) belongs to a four-dimensional Extended Chebyshev System and therefore has at most three distinct real zeros (Sections 29.5--29.6).
\item Consequently, every interior critical point collapses to at most three distinct log-likelihood values, independently of the ambient dimension \(n\) (Theorem~\ref{ch:thm:29-3}).
\item The problem reduces to a finite-dimensional algebraic system (Corollary~\ref{ch:cor:29-2}), explicitly characterized in Section 29.8.
\item The global maximum may lie on the boundary of the simplex, which requires separate analysis (Section 29.9).
\end{itemize}
Theorem~\ref{ch:thm:29-3} characterizes \textbf{interior regular constrained critical points}. No claim is made regarding the existence, uniqueness, or global optimality of such points. Likewise, the possibility that global maximizers occur on the boundary of the simplex is not excluded.
The dimensional collapse phenomenon identified here is a direct consequence of the interaction between the stationarity equations and the Extended Chebyshev structure of the space spanned by \(\{1,x,e^x,xe^x\}\). This interaction is independent of the dimension \(n\) and therefore provides a universal structural result for the variational geometry of the Uniform Observer Divergence.
